\documentclass[oneside]{sn-jnl}
\usepackage[T5,T1]{fontenc}
\usepackage[utf8]{inputenc}
\usepackage[english]{babel}
\usepackage{indentfirst}
\usepackage{graphicx}
\usepackage{amsmath}
\usepackage{booktabs}
\usepackage{tabularx}
\usepackage{xcolor}
\usepackage[numbers,sort&compress]{natbib}
\geometry{a4paper,margin=2.5cm}
\setcitestyle{numbers,square,sort&compress}
\newcommand{\pending}[1]{\textbf{\textcolor{red!65!black}{[PLACEHOLDER: #1]}}}
\newcommand{\artifact}[1]{\nolinkurl{#1}}
\raggedbottom

\begin{document}
\begin{titlepage}
\centering
\begin{minipage}{.45\textwidth}
\IfFileExists{logo.png}{\includegraphics[width=5cm]{logo.png}}{}
\end{minipage}\hfill
\begin{minipage}{.5\textwidth}
\raggedleft\textbf{MINISTRY OF EDUCATION AND TRAINING}
\end{minipage}

\vspace{2cm}
{\Huge\bfseries FPT UNIVERSITY\par}
\vspace{1.8cm}
{\huge Capstone Project Document\par}
\vspace{.6cm}
{\Large\bfseries VaxiS: An Evidence-Grounded Retrieval-Augmented Generation System for Vietnamese Vaccine Information Retrieval\par}
\vspace{1.8cm}

\begin{tabularx}{.95\textwidth}{|p{.26\textwidth}|X|}
\hline
\textbf{Group} & AIP491 Group 1 \\
\hline
\textbf{Group members} &
Nguyễn Mạnh Hùng (HE194092), Trần Trọng Nghĩa (HE190336),
Hồ Minh Hiếu (HE194091). \\
\hline
\textbf{Supervisor} & Trần Văn Hà \\
\hline
\textbf{Capstone project code} & AIP491 \\
\hline
\end{tabularx}
\vfill
Hanoi, \pending{confirm official submission month and year}
\end{titlepage}

\begin{center}\Large\bfseries Abstract\end{center}

Vietnamese vaccine information is published in documents that vary in structure and are often written for specialist readers. VaxiS is a Vietnamese-language retrieval-augmented generation system designed to help users locate relevant passages and trace an answer to its documentary evidence. The repository contains a structure-aware Markdown indexing pipeline, lexical and dense retrieval paths, and a hybrid retrieval pipeline that combines dense and learned sparse rankings before optional reranking. A shared generator uses the retrieved chunks as context and tracks citations to chunk identifiers. The experimental design follows six configurations in order: Naive, Advanced, Hybrid (baseline), Table-aware, Graph RAG, and Agent RAG. The working benchmark contains 300 Vietnamese questions spanning single-hop, multi-hop, and table-reasoning tasks. Saved retrieval-only reports cover all 300 benchmark questions at a retrieval depth of five. Naive, Advanced and Hybrid achieve Hit@5 values of 0.5367, 0.5900 and 0.7433, respectively. Preliminary Table-aware and LightRAG-augmented reports are also available, although their implementation documentation and full evaluation remain pending. Agent RAG evaluation remains pending. Answer-quality, citation-support and abstention results are still needed to assess the full question-answering system.

\noindent\textbf{Keywords.} Vaccine information, Vietnamese question answering, retrieval-augmented generation, hybrid retrieval, evidence attribution.

\newpage
\begin{center}\Large\bfseries Acknowledgments\end{center}
We thank our supervisor, Trần Văn Hà, for supervising this capstone project.
\pending{Confirm any additional acknowledgments and permission to name contributors or reviewers.}

\newpage
\tableofcontents
\newpage
\listoffigures
\newpage

\section{Introduction}

\subsection{Motivation}
VaxiS addresses a practical information access problem. A user may know a vaccine by its product name and have a specific question, while the relevant source is a long instruction leaflet or regulatory document. Locating the applicable passage requires identifying the product, recognizing the relevant section and interpreting any qualifications recorded in nearby text or table notes. The project supports natural-language lookup of information in the collected vaccine usage instructions. The topic taxonomy in the older testdata module belongs to its dataset-generation workflow and does not define the current system's supported question categories.

Retrieval-augmented generation (RAG) brings evidence from an external collection into the generation process~\cite{lewis2020rag}. In this setting, the evidence should be accessible to the reader so the answer can be checked against the source. Producing a fluent response, citing a retrieved passage, and supporting each factual claim are distinct requirements. Research on generation with citations similarly evaluates attribution as a quality dimension alongside answer correctness~\cite{gao2023citations}.

The primary intended users are ordinary citizens. Healthcare workers and healthcare students are secondary users who may use the system for document lookup. The scope is general access to published information. Diagnosis, individual treatment decisions, and processing personal medical records fall outside the project scope.

\subsection{Objectives}
The main research question is how to retrieve relevant evidence from Vietnamese vaccine documents and use it to generate a comprehensible answer with traceable citations, including when the available evidence is incomplete.

\begin{enumerate}
\item Prepare a searchable corpus from available vaccine source documents, retaining the document identifiers and section metadata needed to trace evidence.
\item Compare lexical, dense, and hybrid retrieval, then assess the proposed Table-aware, Graph RAG, and Agent RAG extensions.
\item Generate Vietnamese responses that cite the supplied evidence and communicate when the evidence is insufficient.
\item Assess retrieval quality, answer quality, citation support, and abstention using documented datasets and reproducible run configurations.
\end{enumerate}

\section{Related Work}

\subsection{Automated Access to Vaccine Information}
Li et al.\ developed VaxBot-HPV around a curated knowledge base and a fine-tuned GPT model for HPV vaccine questions. Their evaluation included answer relevancy and faithfulness, two criteria that inform the assessment of VaxiS~\cite{li2025vaxbothpv}.

The Vaccine Information Resource Assistant (VIRA), developed by Johns Hopkins and IBM Research, uses natural language understanding to match questions to expert-vetted responses about COVID-19 vaccination. Its English and Spanish interfaces illustrate the value of organizing answers around user concerns and making reviewed content accessible in the user's language~\cite{jhuvira}.

K-Bot provides culturally tailored HPV vaccine information in Korean and English. Its authors designed dialogs around evidence-based content adapted to their intended audience. This work motivates attention to phrasing and accessibility in a Vietnamese-language interface~\cite{kim2025kbot}. VaxiS focuses on source-grounded question answering over the Vietnamese documents available to the project, with provenance carried through retrieval and generation.

\subsection{Evidence-Grounded Question Answering over Vaccine Documents}
BM25 uses corpus statistics and document length to score lexical matches~\cite{robertson2009bm25}. Dense retrieval instead ranks passages through similarity in an embedding space. BGE-M3 provides both dense representations and learned sparse lexical weights in a multilingual encoder~\cite{chen2024m3}. Reciprocal rank fusion combines the positions assigned by multiple retrieval streams on a common scale~\cite{cormack2009rrf}.

Evidence can also cross representational boundaries. HybridQA studies questions that require combining tables with associated passages~\cite{chen2020hybridqa}. Graph RAG research illustrates how entities and relationships extracted from documents can organize evidence for a query~\cite{edge2024graphrag}. LightRAG combines graph indexing with dual-level retrieval and provides the basis for the graph-augmentation candidate in the current repository~\cite{guo2024lightrag}. ReAct illustrates how a language model can interleave decisions with actions that gather information~\cite{yao2023react}. These works motivate candidate extensions to the VaxiS pipeline. The specific designs and empirical effects of those extensions require project implementation records and completed runs.

\section{Methodology}

\subsection{Overview of the Vaccine Question-Answering System}
The documented pipeline has a document-preparation and indexing stage followed by question answering. Document preparation collects PDFs from the Drug Administration of Vietnam public service portal and converts them to Markdown using Google's Gemini. Indexing turns cleaned Markdown into identifiable chunks, adds document and section metadata, and persists the chunks in Parquet. The vector indexing path prepares dense and sparse representations in Qdrant. During question answering, the selected retriever ranks chunks. The generator receives the question and selected evidence, produces a Vietnamese response, and tracks the chunk identifiers cited by the response. The evaluation tooling also supports scoring ranked results stored from an earlier run.

\begin{figure}[htbp]
\centering
\fbox{\parbox{.9\linewidth}{\centering
Scraped PDFs $\rightarrow$ Gemini $\rightarrow$ Cleaned Markdown\\[.6em]
Section-aware chunks with metadata $\rightarrow$ Parquet and Qdrant\\[.6em]
Vietnamese question $\rightarrow$ Selected retriever $\rightarrow$ Evidence chunks\\[.6em]
Evidence and question $\rightarrow$ Generator $\rightarrow$ Response and source identifiers}}
\caption{Documented VaxiS indexing and question-answering flow. The evidence store used depends on the retrieval configuration.}
\label{fig:system}
\end{figure}

The repository includes a conversational entry point with query planning and memory, as well as single-query evaluation entry points. The benchmark protocol should record the entry point used. The presence of conversational planning in the application is insufficient to establish the Agent RAG experiment.

\subsection{Vaccine Metadata Module}
The corpus organization and vaccine name mapping associate cleaned source records with product identifiers and names. The implemented chunk schema supplies document identity, section headings, text, a topic category, a content type and a table flag. These fields support traceability and subsequent analysis, for example by identifying table-bearing chunks or grouping evidence by product. The document-discovery code concatenates the Markdown files in each registration folder into a single indexing document. It constructs a document ID from the mapped vaccine name and folder ID, and synthesizes a display title beginning with the Vietnamese phrase for vaccine usage instructions. This display title does not establish the title or type of each original PDF. The stored source-file metadata records Markdown filenames. The project plan also identifies document type, issuer, publication date and regulatory status as useful provenance fields.
\pending{Supply the frozen corpus metadata dictionary, field coverage, source URL and original-page mapping policy, and the rule for selecting among conflicting or superseded documents.}

\subsection{Vaccine Document Processing Module}
The team collected vaccine usage instructions from the Drug Administration of Vietnam public service portal at \url{https://dichvucong.dav.gov.vn/} on September 18, 2026. The listing was obtained with an HTTP GET request whose limit parameter was increased to retrieve all 123 vaccine listings in one request. The document extraction code supports several document fields, while the corpus selected for this project contains usage instructions. Some instructions also contain authorization-decision or regulatory-guidance information within the same PDF.

The preparation work was divided into three parts. PDFs were trimmed to retain the information needed for the project, removing label pages and other unnecessary material. Each file was uploaded individually to the Gemini website and converted to Markdown using the website's default model. Table images were also converted to HTML using Gemini, and the HTML was inserted into the Markdown files. The converted files were manually reviewed afterward. The exact Gemini model version was not recorded in the available materials.

The repository indexing code starts from the cleaned Markdown. Its parser splits the content into sections, and its chunker groups short sections and divides long sections under a length budget. The chunker defaults to a minimum of 80 and maximum of 600 estimated token units, with an overlap of 50. The estimate uses whitespace-separated words and an added allowance for HTML tags. The code accounts for Markdown and HTML table markup and handles oversized tables separately from prose. These are the current implementation defaults. \pending{Confirm the chunking settings used to build the corpus evaluated in the experiments.} The distinct Table-aware ablation still requires a versioned definition of its representation and retrieval behavior.

The design must preserve provenance when chunk boundaries are changed. A question can depend on a qualifier in an adjacent paragraph or a footnote that changes the interpretation of a row. A useful audit records the original source location and whether the selected evidence preserves those dependencies.
\pending{Provide the conversion prompt if retained, the conversion dates, and any recorded details of the manual review procedure. The exact Gemini model version remains unspecified unless a contemporaneous record is available.}

\subsection{Vaccine Knowledge Graph Module}
Graph RAG is the fifth configuration in the study. Its purpose is to test whether graph-structured evidence helps retrieve information distributed across passages. The current repository includes \artifact{scripts/lightrag_ingestion.py}, graph storage artifacts, and \artifact{retrieval/lightrag_retriever.py} and \artifact{retrieval/hybrid_lightrag.py} within the project package. The M4 report records LightRAG hybrid mode with reciprocal rank fusion and reranking over the table-aware collection. These artifacts establish a preliminary graph-augmentation path. Its mapping to the final Graph RAG configuration and the complete implementation specification require confirmation. A reproducible specification requires the graph schema, the extraction and entity resolution rules, and the links from graph facts to source evidence.
\pending{Provide the Graph RAG implementation path and commit, node and edge schema, entity normalization and extraction procedure, storage backend, graph coverage and counts, and provenance strategy. Attach a verified graph example.}
\pending{Provide the final graph-specific evaluation slice, complete run manifest and answer-level metrics, and observed failure cases beyond the preliminary retrieval results below.}

\subsection{Retrieval-Augmented Generation Module}
Naive uses the repository's BM25 retrieval path over the chunk corpus. The evaluation CLI selects it with \artifact{--baseline}. Hybrid is the study baseline against which subsequent extensions are compared. Advanced embeds the original query with BGE-M3 and searches the dense Qdrant index before passing the top-ranked chunks to the shared generator.

Hybrid queries the BGE-M3 dense and learned sparse indexes. The implementation aggregates candidate ranks using weighted reciprocal rank fusion,
\[
 S(d)=\sum_{i:\,d\in L_i}\frac{w_i}{c+\mathrm{rank}_i(d)},
\]
where \(L_i\) is a candidate ranking, \(w_i\) its weight and \(c\) a smoothing constant. Hybrid uses the \artifact{BAAI/bge-m3} embedding checkpoint and the \artifact{BAAI/bge-reranker-v2-m3} cross-encoder reranker. The current configuration defaults to 50 candidates from each stream, a fusion constant of 60, 30 fused candidates, a final depth of five and enabled reranking. These are code defaults, and the experiment manifest must establish the settings actually used. The run record must identify the candidate depths, fusion settings, reranker checkpoint and final evidence depth.

Table-aware is the fourth configuration and is intended to measure retrieval and grounding on questions that depend on tabular structure. The current code in \artifact{indexing/table_pipeline.py} and \artifact{indexing/table_summarizer.py} implements a summarize\_dual option within the project package. It merges consecutive table fragments within a document and section, uses a prose summary as embedding text, and retains the summary alongside the original merged table for generation. The metadata retains merged source IDs. The results summary identifies M3b with this option and the table-dual collection. The evaluation report records the collection name but does not record the index-building command or table-summary model, so these run details remain pending.
\pending{Provide the Table-aware implementation path and commit, table and cell representation, preservation of headers and footnotes, retrieval and reranking changes, and source attribution behavior. State exactly how it differs from the shared chunker and Hybrid.}
\pending{Provide the final Table-aware indexing manifest and summary settings, confirm the table-dependent subset labels, and add generated answers, answer-level metrics and an audited row or footnote error example.}

For Graph RAG, the retrieval component must specify how a question is associated with graph entities, how candidate paths are selected, and how their supporting passages enter the generator.
\pending{Provide the actual Graph RAG retrieval algorithm, graph search parameters, evidence ordering, and fallback behavior.}

Agent RAG is the sixth configuration and is intended to evaluate tool-mediated evidence gathering across retrieval steps.
\pending{Provide the Agent RAG implementation path and commit, model, available tools, routing and stopping logic, evidence aggregation and citation rules, call budget, fallback behavior, and an example trace.}
\pending{Provide the Agent RAG evaluation protocol, query coverage, comparison settings, metrics, tool-call and latency summaries, and failure analysis.}

\subsection{Answer Grounding and Verification}
The generator asks the model to answer from supplied evidence, cite chunk identifiers and communicate when the context is insufficient. The implementation extracts citation identifiers and checks their membership in the supplied chunks. It records valid and invalid identifiers and flags refusals by matching a fixed list of Vietnamese phrases. This heuristic requires validation against human judgments. The current shared LLM response-validation function rejects empty or non-text content, explicit provider refusals and unexpected completion finish reasons. The planner applies additional output checks. Validation behavior should be tied to the code revision used for a reported run. Citation membership establishes traceability to context. Assessing whether a passage supports a particular claim requires semantic review.

The separate review application gives reviewers access to questions, reference answers, and source passages. Its workflow supports changes to the answer and evidence selection, with code paths for consensus and adjudication. The review application provides a way to audit the benchmark and its provenance.
\pending{Provide the final expert review protocol, reviewer qualifications and counts, adjudication outcomes, citation support rubric, and whether any separate automated claim-verification module contributed to the reported runs.}

\section{Dataset}
The corpus contains 75 vaccine usage-instruction PDFs collected from the Drug Administration of Vietnam public service portal on September 18, 2026. The team retrieved 123 vaccine listings. Of these, 44 had no documents, one could not be downloaded because its URL was malformed, two produced empty folders, and one contained only a label with no usage instructions. Excluding these 48 listings left 75 PDFs. Some included PDFs also contain authorization-decision or regulatory-guidance information.

The team confirmed that the 75 PDFs correspond to the 75 document IDs in the current corpus. Inspection of the checked-in Parquet artifact finds 1,744 unique chunks, including 490 with the table flag set. The current Qdrant store is used for retrieval. \pending{Confirm whether that store was built from the checked-in 1,744-chunk Parquet artifact or an earlier version of it.}

The main 300-question benchmark is generated using DeepEval's Synthesizer. In \artifact{scripts/generate_eval_dataset.py}, prepared chunk contexts are passed to \artifact{generate_goldens_from_contexts}, with \artifact{include_expected_output=True} and one golden requested per context group. Styling and evolution configurations control answerability, question style and reasoning complexity. DeepEval documents this interface as generating questions and expected answers from supplied contexts~\cite{deepevalcontexts}. The script constructs context pools for single-hop, multi-hop and table-reasoning questions, including multi-hop evidence from the same document or different vaccines. The saved benchmark metadata records \artifact{gpt-oss-120b} as the generation model. \pending{Confirm the final benchmark export, generation command, script revision, DeepEval version and model endpoint.}

The older \artifact{testdata} module and its question-category taxonomy are legacy artifacts. The team uses the older query dataset only for unanswerable-query testing. It does not supply the main benchmark or define the current system's question categories. The checked-in \artifact{test_dataset.json} contains 460 records with 74 answerable, 164 partially answerable and 222 unanswerable labels. Its additional evaluation pass currently loads the entire file unless a filtered artifact is supplied, so the report must identify the subset actually scored for unanswerable-query testing. \pending{Confirm which file or exported subset is used for unanswerable testing, how its labels were checked and its final scored count. The team referred to test\_data.json; the checked-in filename is test\_dataset.json.}

The 300-question file stores a question, expected output, context, and metadata containing evidence identifiers and question attributes. Table~\ref{tab:dataset} gives the label counts observed in \artifact{evaluation/eval_dataset_300_final_fixed.json}. Each dimension describes the same set of 300 records.

\begin{table}[htbp]
\centering
\caption{Label distribution in the 300-question working benchmark.}
\label{tab:dataset}
\begin{tabular}{llr}
\toprule
Dimension & Label & Count\\
\midrule
Answerability & Answerable & 210\\
 & Partially answerable & 90\\
Complexity & Single-hop & 150\\
 & Multi-hop & 90\\
 & Table reasoning & 60\\
Question style & Patient style & 180\\
 & Clinical style & 120\\
\bottomrule
\end{tabular}
\end{table}

The corrected 300-question file changes ordered source-chunk lists for 279 records relative to the earlier file. Its goldens match those embedded in the saved generation outputs. This establishes which labels accompany the saved outputs. Expert validation of label completeness and answer support must be documented separately. The 300-question file contains answerable and partially answerable labels, so evaluation of fully unanswerable questions requires a declared additional sample.
\pending{Provide the final adjudicated dataset export, version checksum, sampling and split procedure, overlap checks between artifacts, and the size and selection rule for any fully unanswerable subset.}

\section{Experiments and Results}

\subsection{Experimental Design}
The study compares the six configurations in Table~\ref{tab:sequence}. The newly available baseline reports record live retrieval on \artifact{evaluation/eval_dataset_300_final_fixed.json} with top-k five and cutoffs of one, three and five. Each scores all 300 questions, with no short rankings. These are deterministic chunk-retrieval evaluations. Their empty generated-answer fields provide no observations of answer quality or refusal behavior. Each comparison should use a frozen corpus and judgment set, the same scored questions and matched evidence budgets. Record differences intrinsic to the methods, such as using a graph traversal or multiple retrieval calls. Hold the generator, prompts and judge settings fixed for comparisons intended to isolate the effect of retrieval. Hybrid serves as the reference for extensions. Because Hybrid can change both retrieval streams and reranking, the sequence is a system-level comparison. Component effects require matched controls, such as Hybrid with reranking disabled.

\begin{table}[htbp]
\centering
\caption{Ablation-study sequence and documented evidence paths.}
\label{tab:sequence}
\begin{tabularx}{\linewidth}{rlX}
\toprule
Order & Configuration & Evidence path\\
\midrule
1 & Naive & BM25 over corpus chunks.\\
2 & Advanced & One BGE-M3 dense query against Qdrant.\\
3 & Hybrid (baseline) & Dense and learned sparse retrieval, rank fusion and optional reranking.\\
4 & Table-aware & \pending{Record final table representation and evidence path.}\\
5 & Graph RAG & \pending{Record final graph construction and retrieval path.}\\
6 & Agent RAG & \pending{Record final orchestration and tool-use path.}\\
\bottomrule
\end{tabularx}
\end{table}

For retrieval, Hit@$k$ indicates whether at least one labeled relevant chunk was retrieved by rank \(k\). Recall@$k$ measures the proportion of labeled evidence recovered and Precision@$k$ measures the proportion of the first \(k\) positions occupied by labeled evidence. Mean reciprocal rank rewards early retrieval of the first labeled source. Normalized discounted cumulative gain rewards placing labeled sources near the top. Alternative supporting chunks can receive no credit when the gold evidence is incomplete. Report scoring coverage and the number of rankings shorter than each cutoff.

The evaluation code also supports mean average precision at a cutoff and LLM-judged answer relevancy and faithfulness, contextual relevancy, precision and recall, and a default-enabled GEval Correctness metric against the reference answer. Citation analysis should distinguish citation validity from claim support, and failure analysis should assess whether refusals and partial answers match the evidence available. The source ODT identifies user testing as an evaluation objective.
\pending{Provide the final run manifest with code revision, index revision, checkpoint identifiers, prompts, judge settings, hardware, retrieval and context budgets, inclusion and exclusion rules, and significance or uncertainty procedure if used. Supply the human review and user-testing protocols and completed results if available.}

\subsection{Results}
The implementation audit for this report update inspected code commit \artifact{7798195}. The saved run reports do not record their generating code commits. The principal baseline sources are \artifact{evaluation/ir_naive_bm25.json}, \artifact{evaluation/ir_advanced.json} and \artifact{evaluation/ir_hybrid.json}. Inspection of their query text, reference answers and labeled evidence IDs finds agreement with the current corrected benchmark for all 300 records. Their stored per-query scores reproduce with the current deterministic scorer, and their overall means agree with those scores to the four decimal places reported. This is a static audit of saved outputs, with no new retrieval run.

\begin{table}[htbp]
\centering
\caption{Saved retrieval-only baseline results on 300 questions. Each run retrieves five chunks.}
\label{tab:results}
\begin{tabular}{lrrr}
\toprule
Metric & Naive & Advanced & Hybrid (baseline)\\
\midrule
Hit@1 & 0.2767 & 0.3067 & 0.4233\\
Hit@3 & 0.4533 & 0.4767 & 0.6767\\
Hit@5 & 0.5367 & 0.5900 & 0.7433\\
Recall@5 & 0.4733 & 0.5167 & 0.6817\\
Precision@5 & 0.1153 & 0.1280 & 0.1680\\
MRR@5 & 0.3747 & 0.4064 & 0.5532\\
NDCG@5 & 0.3773 & 0.4093 & 0.5607\\
MAP@5 & 0.3292 & 0.3562 & 0.5019\\
\bottomrule
\end{tabular}
\end{table}

Naive retrieves at least one labeled source in the first five positions for 161 questions, Advanced for 177, and Hybrid for 223. Hybrid therefore has 62 more successful questions than Naive and 46 more than Advanced, corresponding to net differences of 20.67 and 15.33 percentage points in Hit@5. Advanced exceeds Naive by 16 questions, or 5.33 percentage points. These aggregate comparisons describe the saved configurations. Isolating the contribution of sparse retrieval from reranking requires additional controls. No significance tests or repeated-run variability estimates are included in these artifacts.

\begin{table}[htbp]
\centering
\caption{Baseline Hit@5 and Recall@5 by benchmark complexity.}
\label{tab:slices}
\begin{tabular}{lrrrrrr}
\toprule
 & \multicolumn{2}{c}{Naive} & \multicolumn{2}{c}{Advanced} & \multicolumn{2}{c}{Hybrid}\\
Group & Hit@5 & Recall@5 & Hit@5 & Recall@5 & Hit@5 & Recall@5\\
\midrule
Single-hop (150) & 0.6200 & 0.6200 & 0.6267 & 0.6267 & 0.8133 & 0.8133\\
Multi-hop (90) & 0.5556 & 0.3444 & 0.6556 & 0.4111 & 0.7333 & 0.5278\\
Table reasoning (60) & 0.3000 & 0.3000 & 0.4000 & 0.4000 & 0.5833 & 0.5833\\
\bottomrule
\end{tabular}
\end{table}

For multi-hop questions, Hybrid Hit@5 is 0.7333 while Recall@5 is 0.5278. Finding one labeled source leaves part of the labeled evidence unrecovered in some cases. The table-reasoning subset has the lowest Hit@5 for each baseline, motivating a closer assessment of table representation. For partially answerable questions, Hit@5 is 0.3667, 0.4556 and 0.6667 for Naive, Advanced and Hybrid. These figures assess retrieval against labels and leave the appropriateness of the generated partial answer unresolved.

\begin{table}[htbp]
\centering
\caption{Preliminary extension retrieval results. M3b and M4 use merged-source alias matching. Complete implementation and answer-quality evaluation remain pending.}
\label{tab:extensions}
\begin{tabular}{lrrrr}
\toprule
Metric & Hybrid M2 & M3b strict & M3b alias & M4 LightRAG\\
\midrule
Hit@1 & 0.4233 & 0.4567 & 0.4933 & 0.4933\\
Hit@3 & 0.6767 & 0.6667 & 0.7300 & 0.7300\\
Hit@5 & 0.7433 & 0.7233 & 0.7833 & 0.7833\\
Recall@5 & 0.6817 & 0.6533 & 0.7167 & 0.7167\\
Precision@5 & 0.1680 & 0.1613 & 0.1753 & 0.1753\\
MRR@5 & 0.5532 & 0.5643 & 0.6114 & 0.6108\\
NDCG@5 & 0.5607 & 0.5597 & 0.6122 & 0.6115\\
MAP@5 & 0.5019 & 0.5089 & 0.5581 & 0.5570\\
\bottomrule
\end{tabular}
\end{table}

The extension sources are \artifact{eval_ir_m2_hybrid.json} for M2, \artifact{eval_ir_m3b_table_dual_before_preserving.json} for M3b strict, \artifact{eval_ir_m3b_table_dual.json} for M3b alias, and \artifact{eval_ir_m4_hybrid_lightrag_rrf.json} for M4. The M3b strict and alias-aware files contain the same ranked IDs for every question. Alias-aware scoring credits the source IDs covered by a merged table chunk while preserving its single ranking position. Alias mappings are recorded for 115 questions. At depth five, the score change recovers 18 previously uncredited hits, moving Hit@5 from 217/300 to 235/300. This six-percentage-point difference comes entirely from the scoring rule. Compared with Hybrid M2, M3b alias has 12 additional successful questions overall, or four percentage points. The contrast includes a different evidence representation and its associated matching rule. Under the current alias scorer, Precision counts ranking slots that cover new gold evidence, Recall counts covered source IDs, and MAP and NDCG adjust their relevant-item denominators for multiple sources represented in a single slot. These definitions must accompany the extension scores.

Within the 60 table-reasoning questions, M2 retrieves a labeled source for 35 questions, M3b strict for 29 and M3b alias for 46. The M3b alias difference from M2 is 11 questions, or 18.33 percentage points. Its difference from M3b strict is 17 questions, or 28.33 percentage points, attributable to alias credit on unchanged rankings. For single-hop questions, M3b alias Hit@5 decreases from Hybrid's 0.8133 to 0.8067. Improvement is therefore uneven across slices.

M4 records \artifact{hybrid_lightrag} retrieval with \artifact{rrf_rerank} fusion. It shares M3b alias's aggregate Hit, Recall and Precision at the reported cutoffs, with slightly lower MRR@5 and NDCG@5. Its top-five order differs from M3b for nine questions, and its top-five membership differs for six. Identical aggregate coverage does not establish identical per-question evidence. For 297 questions, the saved ranking contains five document IDs plus a sixth \artifact{lightrag_kg_context} ID. The graph-context entry lies outside the evaluated top five. These scores leave any answer-generation benefit from that context unmeasured.

\pending{Provide final Table-aware indexing and summary settings, verify the intended mapping of M4 to the final Graph RAG configuration, and supply the complete implementation records and answer-quality evaluations for both extensions. Provide Agent RAG implementation and evaluation results.}

The earlier \artifact{evaluation/eval_dataset_300_generation_results.json} contains four retrieved IDs per response and is a separate generation artifact. It cannot identify the generator or answer quality of these newer top-five IR runs. \pending{Provide corresponding generated answers, judge configuration, faithfulness, correctness, citation-support, unanswerable-query and latency results. Record code, corpus and index revisions and the final checkpoint identities for each run.}

\subsection{Qualitative Assessment of Vietnamese Vaccine Answers}
A qualitative assessment should trace the question to the ranked evidence and examine the generated response claim by claim. For product-specific questions, check whether the passage refers to the right product and formulation. For tables, check whether the evidence includes the relevant headers and any qualifying notes. For evidence distributed across sources, check whether all required passages appear and whether the answer combines them consistently. A partially answerable case should show which parts the source supports and how the response handles the remaining uncertainty.

The saved rankings supply two concrete retrieval examples. In \artifact{eval_0003}, a Nimenrix question has labeled evidence \artifact{nimenrix_540310196925_c_0000}. Naive does not recover it in five results, while Advanced and Hybrid place it second. In \artifact{eval_0170}, a Vaxneuvance table question has labeled evidence \artifact{vaxneuvance_539310049125_c_0054}. Hybrid misses it at depth five, Advanced places it third, and preliminary M3b places it second. These examples show query-specific ranking differences. Human review of the source passages and generated answers is still required for judgments about clinical correctness or answer support.

\pending{Insert reviewed examples with query IDs, variant, source titles and chunk IDs, source excerpts, actual output and assessor judgments. Include a supported answer, an evidence miss or unsupported claim, a table-dependent example and a partially answerable case. Supply reviewed Graph RAG and Agent RAG examples once the corresponding runs are available.}

\section{Conclusion}
VaxiS applies evidence-grounded retrieval and generation to Vietnamese vaccine information. The available project materials establish an indexing path from cleaned Markdown, BM25, dense and hybrid retrieval implementations, chunk-level citation tracking, a working benchmark, and tooling for reviewing reference answers and evidence. These artifacts support evaluation of whether relevant passages are retrieved and whether the resulting answers can be checked against those passages.

The saved baseline evaluations show higher chunk-retrieval coverage for Hybrid than for Naive or Advanced on the 300-question benchmark at depth five. Table reasoning remains the weakest baseline slice. Preliminary Table-aware results show the importance of preserving merged-source identities during evaluation, and the available graph-augmentation report shows no aggregate coverage increase over the alias-aware Table-aware configuration at the reported cutoffs. Full implementation documentation for these extensions, Agent RAG evaluation and answer-level assessment remain pending. These retrieval findings do not establish clinical correctness, improved generation or user benefit.
\pending{Complete the conclusion with verified answer-level, citation-support, abstention and user-evaluation findings when available, and retain the limits of the evaluated corpus and benchmark.}

\bibliographystyle{unsrtnat}
\bibliography{references}

\appendix
\section*{Appendix}
\addcontentsline{toc}{section}{Appendix}

\subsection{Source-to-Markdown Processing, Metadata Extraction, and Chunking}
Collection is documented in \artifact{scrape_dataset/scrape_data.py}, \artifact{scrape_dataset/extract_documents.py}, and \artifact{scrape_dataset/download_documents.py}. The team used the Gemini website with its default model, uploading one trimmed PDF at a time. Table images were converted to HTML and inserted into the Markdown, followed by manual review. \pending{Provide the conversion prompt or retained session record if available.}

The principal indexing files are \artifact{src/vacxis_capstone_fa26/indexing/parser.py}, \artifact{src/vacxis_capstone_fa26/indexing/chunker.py}, and \artifact{src/vacxis_capstone_fa26/indexing/indexer.py}. Review the original source alongside its cleaned Markdown and indexed output when checking the preservation of a table or footnote.
\pending{Provide a representative source-to-Markdown example, Gemini conversion and cleaning settings, and the indexing command and resulting chunk metadata for the final corpus.}

\subsection{Answer Generation Prompts and Output Schema}
The generation files \artifact{src/vacxis_capstone_fa26/generation/prompts.py} and \artifact{src/vacxis_capstone_fa26/generation/generator.py} define the prompt construction and citation handling. Saved result records include the answer, valid and invalid citations, cited and retrieved IDs, a refusal flag, model, latency and usage.
\pending{Insert the exact versioned generator prompt, context formatting, model settings and a sanitized saved response record from the final comparison.}

\subsection{Graph Schema and Constraints}
\pending{Insert the versioned Graph RAG node and edge specification, graph serialization or storage constraints, provenance fields, and a worked retrieval trace with source identifiers.}

\subsection{Example Questions and Reviewed Outputs}
\pending{Insert approved benchmark examples with expected and actual outputs, evidence IDs, and reviewer notes. Include each evaluated configuration where comparable examples are available.}

\subsection{Project Management Plan}
The ODT records a 14-week plan organized around problem definition, preparation of data, development of the AI pipeline, system development and evaluation. Table~\ref{tab:timeline} summarizes the planned milestones.

\begin{table}[htbp]
\centering
\caption{Planned milestones documented in the project ODT.}
\label{tab:timeline}
\begin{tabularx}{\linewidth}{lX}
\toprule
Weeks & Planned activities\\
\midrule
1--3 & Define the problem, scope and research questions. Set milestones and survey related systems.\\
4--6 & Collect and clean data, prepare metadata and chunks, and implement retrieval.\\
7--8 & Integrate generation, citations and evidence handling. Refine prompts and evaluation.\\
9--11 & Develop the backend and chatbot interface and integrate the end-to-end system.\\
12--14 & Run evaluation, analyze errors, assemble results and complete the report.\\
\bottomrule
\end{tabularx}
\end{table}

The plan identifies OCR or metadata errors, retrieval misses, unsupported generated claims, and the difficulty of assessing answer and citation quality as project risks. Its proposed mitigations include checking processed documents against originals, comparing retrieval strategies, constraining generation to source evidence, and stratifying evaluation by question type.
\pending{Confirm the actual milestone dates, deliverables and deviations from the plan.}
\end{document}
